\subsubsection{c}

Akselerasjonen i vertikal retning er gitt ved:
\begin{align}
a(t) = kt
\end{align}

Vi finner først farten ved å integrere akselerasjonen:
\begin{align}
v(t) = \int a(t)\,dt = \int kt\,dt = \frac{1}{2}kt^2 + C
\end{align}

Siden kula starter uten vertikal starthastighet, har vi $v(0)=0$, som gir $C=0$:
\begin{align}
v(t) = \frac{1}{2}kt^2
\end{align}

Videre finner vi posisjonen:
\begin{align}
s(t) = \int v(t)\,dt = \int \frac{1}{2}kt^2\,dt = \frac{1}{6}kt^3 + C
\end{align}

Med startbetingelsen $s(0)=0$ får vi:
\begin{align}
s(t) = \frac{1}{6}kt^3
\end{align}

Kula treffer den nederste platen når $s(t) = 0.40\,\text{m}$:
\begin{align}
\frac{1}{6}kt^3 = 0.40
\end{align}

Løser for $t$:
\begin{align}
t = \sqrt[3]{\frac{2.4}{k}}
\end{align}

Setter inn $k = 3.00\,\text{m/s}^3$:
\begin{align}
t = \sqrt[3]{\frac{2.4}{3.00}} \approx 0.928\,\text{s}
\end{align}

I horisontal retning er farten konstant:
\begin{align}
v_x = 0.50\,\text{m/s}
\end{align}

Strekningen blir da:
\begin{align}
s_x = v_x t = (0.50\,\text{m/s})(0.928\,\text{s}) = 0.464\,\text{m}
\end{align}